% =========================
% Básicos
% =========================
\usepackage[utf8]{inputenc}
\usepackage[T1]{fontenc}
\usepackage{textcomp}
\usepackage[spanish]{babel}
\usepackage{microtype}
\usepackage{url}
\usepackage{ragged2e}
\usepackage{graphicx}
\usepackage{float}
\usepackage{booktabs}
\usepackage{enumitem}
\usepackage{emptypage}
\usepackage{booktabs} 
\usepackage{tabularx}
\usepackage{subcaption}
\usepackage{multicol}
\usepackage[usenames,dvipsnames]{xcolor}
% \usepackage{hyperref}
% \hypersetup{
%     colorlinks,
%     linkcolor={black},
%     citecolor={black},
%     urlcolor={blue!80!black}
% }

% =========================
% Matemáticas
% =========================
\usepackage{amsmath, amsfonts, mathtools, amsthm, amssymb}
\usepackage{mathrsfs}
\usepackage{cancel}
\usepackage{bm}
\usepackage{systeme}
\usepackage{stmaryrd}

\let\implies\Rightarrow
\let\impliedby\Leftarrow
\let\iff\Leftrightarrow
\let\epsilon\varepsilon

% =========================
% Teoremas y Cajas 
% =========================

% Colores Personalizados
\definecolor{theorembg}{RGB}{245, 255, 250}   
\definecolor{theoremlines}{RGB}{0, 100, 0}   
\definecolor{exercisebg}{RGB}{245, 245, 255}
\definecolor{exerciselines}{RGB}{0, 0, 100}
\definecolor{proofit}{RGB}{50, 50, 50}    
\usepackage{thmtools}
\usepackage[framemethod=TikZ]{mdframed}

\mdfsetup{skipabove=1em, skipbelow=1em}

\declaretheoremstyle[
    headfont=\bfseries\sffamily,
    bodyfont=\normalfont,
    mdframed={
        backgroundcolor=theorembg,
        linecolor=theoremlines,
        linewidth=0.5pt,
        innertopmargin=8pt,
        innerbottommargin=8pt,
        nobreak=true,
    }
]{THMs}

\declaretheoremstyle[
    headfont=\bfseries\sffamily,
    bodyfont=\normalfont,
    mdframed={
        backgroundcolor=exercisebg,
        linecolor=exerciselines,
        linewidth=0.5pt,
        innertopmargin=8pt,
        innerbottommargin=8pt,
        nobreak=true,
    }
]{DEF}

\declaretheoremstyle[
    headfont=\bfseries\sffamily,
    bodyfont=\normalfont,
    mdframed={
        linewidth=1pt,
        rightline=false, topline=false, bottomline=false,
        linecolor=gray!40,
        innerleftmargin=10pt,
    }
]{MSCL}

\declaretheoremstyle[
    headfont=\bfseries\itshape\sffamily,
    bodyfont=\normalfont,
    numbered=no,
    mdframed={
        topline=false, bottomline=false, rightline=false, leftline=false,
    },
    qed=\qedsymbol
]{PROOF}

\declaretheorem[style=DEF, name=Definición]{definicion}
\declaretheorem[style=DEF, name=Axioma]{axioma}

\declaretheorem[style=THMs, name=Teorema]{teorema}
\declaretheorem[style=THMs, name=Proposición]{proposicion}
\declaretheorem[style=THMs, name=Lema]{lema}
\declaretheorem[style=THMs, numbered=no, name=Corolario]{corolario}
\declaretheorem[style=THMs, name=Tarea]{tarea}
\declaretheorem[style=THMs, name=Método]{metodo}

\declaretheorem[style=MSCL, numbered=no, name=Ejemplo]{ejemplo}
\declaretheorem[style=MSCL, numbered=no, name=Observación]{observacion}
\declaretheorem[style=MSCL, numbered=no, name=Nota]{nota}
\declaretheorem[style=MSCL, numbered=no, name=Ejercicio]{ejercicio}
\declaretheorem[style=MSCL, numbered=no, name=Borrador]{borrador}
\declaretheorem[style=MSCL, numbered=no, name=Solución]{solucion}

\let\proof\relax
\let\endproof\relax
\declaretheorem[style=PROOF, name=Demostración]{proof}

% =========================
% Clase...
% =========================
\newcommand{\clase}[3]{%
    \ifthenelse{\isempty{#3}}{%
        \def\@clase{Clase #1}%
    }{%
        \def\@clase{Clase #1: #3}%
    }%
    \subsection*{\@clase}
    \marginpar{\small\textsf{#2}}
}

% =========================
% Cancel Color
% =========================

\newcommand{\colorcancel}[2]{\renewcommand{\CancelColor}{\color{#2}}\cancel{#1}}

% =========================
% Encabezados
% =========================

% fancy headers
\usepackage{fancyhdr}
\pagestyle{fancy}

% \fancyhead[LE,RO]{Gilles Castel}
\fancyhead[RO,LE]{\@clase}
\fancyhead[RE,LO]{}
\fancyfoot[LE,RO]{\thepage}
\fancyfoot[C]{\leftmark}

\makeatother

% =========================
% Figuras externas
% =========================
\usepackage{import}
\usepackage{xifthen}
\usepackage{pdfpages}
\usepackage{transparent}

\newcommand{\incfig}[1]{%
    \def\svgwidth{\columnwidth}
    \import{./figures/}{#1.pdf_tex}
}

\pdfsuppresswarningpagegroup=1

% =========================
% Autor
% =========================
\author{César Pérez Amador}
